\section{Plan de Capacitación y Gestión del Cambio}

\subsection{Objetivo de la gestión del cambio}

La estrategia de capacitación tuvo como propósito principal \textbf{facilitar la transición} del equipo hacia el uso diario del sistema Odoo 19, reduciendo la resistencia al cambio y demostrando el valor del CRM y del tablero de proyectos en la operación real.

\subsection{Metodología de transferencia de conocimientos}

Se desarrollaron las siguientes acciones:

\begin{itemize}
    \item Talleres funcionales para usuarios clave, cubriendo el flujo \textit{de oportunidad a cierre} dentro del CRM.
    \item Sesiones prácticas sobre el uso de la ficha del cliente, el registro de actividades y la gestión del calendario.
    \item Entrega de \textbf{videos tutoriales operativos} como material de consulta permanente para reforzar el aprendizaje.
\end{itemize}

\subsection{Despliegue y Retroalimentación}

El cronograma de capacitaciones se integró con el avance de los sprints de configuración. Esto permitió que la dirección y los usuarios clave validaran la \textbf{usabilidad} del sistema en tiempo real, facilitando ajustes preventivos antes de consolidar el uso en el día a día.

\subsection{Resultados de la adopción tecnológica}

Entre los resultados observados destacan:

\begin{itemize}
    \item Mayor \textbf{autonomía del equipo} en la gestión de contactos y reservas dentro del sistema.
    \item Disminución de la dependencia exclusiva de las conversaciones de \textit{WhatsApp} como fuente de información.
    \item Uso creciente del calendario y de las actividades como soporte al seguimiento diario de los itinerarios.
\end{itemize}

